\documentclass[10pt,a4paper]{amsart}
\usepackage[margin=19mm]{geometry}
\usepackage{amsmath,amssymb,mathtools}
\usepackage[scr=rsfs,cal=euler]{mathalfa}
\usepackage{xfrac}
\usepackage[unicode,hidelinks,bookmarks=false]{hyperref}
\newcommand{\ie}{i.e.\ }
\newcommand{\cf}{cf.\ }
\newcommand{\resp}{respectively\ }
\newcommand{\Subsection}{\subsection}
\providecommand{\nobreakdash}{\nobreak\hbox{-}}
\setlength{\emergencystretch}{2em}
\multlinegap0pt
\usepackage{bm}
\usepackage{bbm}
\usepackage{dsfont}
\usepackage{xspace}
\usepackage{cancel}
\usepackage{mathtools}
\usepackage{cleveref}
\usepackage{amsthm}
\makeatletter
\@ifundefined{theorem}{\newtheorem{theorem}{Theorem}}{}
\@ifundefined{claim}{\newtheorem{claim}[theorem]{Claim}}{}
\@ifundefined{claimproof}{\newtheorem{claimproof}{Claimproof}}{}
\@ifundefined{lemma}{\newtheorem{lemma}[theorem]{Lemma}}{}
\makeatother
\makeatletter
\newcommand\N{\mathbb{N}}
\newcommand\R{\mathbb{R}}
\newcommand{\abs}[1]{\mathopen{}\left\lvert#1\right\rvert\mathclose{}}
\newcommand{\bw}{\overline \omega}
\newcommand{\cA}{\mathcal A}
\expandafter\ifx\csname claimproof\endcsname\relax
\newenvironment{claimproof}[1][Proof]{
	\renewcommand{\oldqed}{\qedsymbol}
	\renewcommand{\qedsymbol}{\endofClaim}
	\begin{proof}[#1]
	}{
	\end{proof}
	\renewcommand{\qedsymbol}{\oldqed}
}
\fi
\newcommand{\coleq}{\coloneq}
\newcommand{\f}{\varphi}
\newcommand{\sm}{\setminus}
\newcommand{\subs}{\subseteq}
\newcommand{\w}{\f\vert_S}
\makeatother
\title[No extremal square-free words over alphabets of size at leas]{No extremal square-free words over alphabets of size at least 5}
\author{Eng Keat Hng, Silas Rathke}
\date{Source: arXiv:2608.10591v1 (CC BY 4.0)}
\begin{document}
\maketitle

\noindent\emph{Source and licence.} No extremal square-free words over alphabets of size at least 5. Version arXiv:2608.10591v1 (CC BY 4.0), distributed under the Creative Commons Attribution 4.0 International licence, as recorded in the arXiv licence metadata for this exact version and shown on the abstract page https://arxiv.org/abs/2608.10591v1. Full rights notice: licenses/arxiv-2608.10591-CC-BY-4.0.txt.\par\smallskip

\noindent\textit{Editorial scope.} Counted unit: the Lemma whose source label is \texttt{lem:bound-same-parity-twothirds} in the article (source0.tex lines 325--328, source label \texttt{lem:bound-same-parity-twothirds}), delivered together with the article's own complete original proof of it (source0.tex lines 330--387, one \texttt{proof} environment of 58 lines). No mathematics is written, repaired or re-derived: statement and proof are the article's own text line for line. Required context retained uncounted: lem:gap-lower-bound (lines 269--281); lem:bound-min-mean-weight-cycle (lines 297--300). Every definition, display and label of those lines is retained unchanged. Omitted from this package and not redistributed: 1--268, 282--296, 301--324, 329--329, 388--637. Nothing inside a retained excerpt is omitted or altered: every retained body is delivered exactly as the article's lines are written, so no omission receipt of any kind accompanies this package. Citations: the retained excerpts carry no citation command at all, so no bibliography entry of the article is transcribed and no reference is redistributed. Reference adaptation: none was needed and none was made; every reference inside the delivered unit resolves inside it, so no unresolved reference or citation is produced. Printed slips kept literal: no printed word, symbol, hypothesis or number of the counted statement or of its proof was altered. Layout and class: the article's own class is \texttt{article} with packages including inputenc, fontenc, comment, adjustbox, mathrsfs, marvosym, dsfont, rotating, amssymb, amsmath, amsthm, thmtools, hyperref, cleveref; the wrapper is the lane's pinned \texttt{amsart} editorial wrapper at 10pt on A4, which restates the theorem-like environments the retained text needs and adds the article's own macro definitions that the retained text needs (\texttt{\textbackslash N}, \texttt{\textbackslash R}, \texttt{\textbackslash abs}, \texttt{\textbackslash bw}, \texttt{\textbackslash cA}, \texttt{\textbackslash coleq}, \texttt{\textbackslash f}, \texttt{\textbackslash sm}, \texttt{\textbackslash subs}, \texttt{\textbackslash w}), with extra wrapper packages (bm, bbm, dsfont, xspace, cancel, mathtools, cleveref). The delivered numbering is the unit's own. Assets: no retained span contains a figure environment, a graphics-inclusion command, a PDF-page inclusion, a table or a TikZ picture; no image file is copied into this package, the package declares no asset dependency and no image file is redistributed with it. Licence: version arXiv:2608.10591v1 of this article is distributed under the Creative Commons Attribution 4.0 International licence, as recorded in the arXiv licence metadata for this exact version and shown on the abstract page https://arxiv.org/abs/2608.10591v1. A full rights notice is delivered at licenses/arxiv-2608.10591-CC-BY-4.0.txt. Characters: every retained excerpt is pure ASCII, so the delivered engine, XeLaTeX with its default Latin Modern text font, reports no missing character.
	\begin{lemma} \label{lem:gap-lower-bound}
		For all~$\ell,\ell'\in\N\sm\{1\}$ we have
		\begin{equation*}
			\f(\ell,\ell') \ge
			\begin{cases}
				\ell &\textrm{if } \ell'=\ell \\
				4\ell-3\ell'+2 &\textrm{if } \ell+1<\ell'\le\frac{4}{3}\ell \\
				1 &\textrm{if } \frac{4}{3}\ell<\ell'<\frac{3}{2}\ell \\
				\ell' &\textrm{if } \frac{2}{3}\ell<\ell'<\ell-1 \\
				0 &\textrm{otherwise.}
			\end{cases}
		\end{equation*}
	\end{lemma}

	\begin{lemma} \label{lem:bound-min-mean-weight-cycle}
		Let~$S\subs\N\sm\{1\}$ be a finite set. Let~$\bw$ denote the minimum mean weight of a directed cycle in the weighted digraph $(G(S),\w)$. Suppose that~$\bw\ne0$. Then for every square-free word~$W$ of length~$n$ we have
		\[ \sum_{\ell\in S}\abs{\cA^{+}_\ell(W)}\le \max\!\left(0,\frac{n+\max S-2\min S+1}{\bw}\right)\,. \]
	\end{lemma}

	\begin{lemma}\label{lem:bound-same-parity-twothirds}
		Let~$W$ be a square-free word of length~$n$. Suppose that~$a,b\in\N\setminus\{1\}$ have the same parity and satisfy~$a\le b<\frac{3}{2}a$. Then
		\[\sum_{i=0}^{\frac{b-a}{2}}\abs{\cA^{+}_{a+2i}(W)}\le \max\!\left(0,\frac{2(n+b+1-2a)}{a+1}\right)\,.\]
	\end{lemma}

	\begin{proof}
		Let $S\coleq\{a,a+2,a+4,\dots,b\}$. By \Cref{lem:bound-min-mean-weight-cycle}, it is enough to show that the minimum mean weight of $(G(S),\w)$ is at least $\frac{a+1}{2}$. To do this, we define the function~$\pi\colon V(G(S))\to\R$ as follows.
		\begin{equation*}
			\pi(u)\coleq
			\begin{cases}
				\frac{a+1}{2}-1 &\textrm{if }u<\frac{3}{4}b \\
				\frac{a+1}{2}-(4u-3b+2) &\textrm{if }\frac{3}{4}b\le u< \frac{a-3+6b}{8} \\
				0 &\textrm{if }\frac{a-3+6b}{8}\le u\le b\,.
			\end{cases}
		\end{equation*}
		We prove the following key claim whose proof entails detailed but elementary calculations with extensive case distinction.
		
		\begin{claim} \label{claim:mean-weight-bound-vtx-pot}
			For all~$(\ell,\ell')\in E(G(S))$ we have~$\w^{\pi}(\ell,\ell')\coleq\w(\ell,\ell')+\pi(\ell)-\pi(\ell')\ge\frac{a+1}{2}$.
		\end{claim}
		Before proving the claim, we will quickly see how it implies that the minimum mean weight of $(G(S),\w)$ is at least $\frac{a+1}{2}$. Let~$C=v_1\dots v_{\ell}v_1$ be a directed cycle in~$G(S)$; set~$v_{\ell+1}\coleq v_1$ for convenience. Then, 
		\[ \frac{1}{\ell}\sum_{i=1}^{\ell}\w(v_i,v_{i+1}) = \frac{1}{\ell}\sum_{i=1}^{\ell}\left(\w(v_i,v_{i+1})+\pi(v_i)-\pi(v_{i+1})\right) \overset{\ref{claim:mean-weight-bound-vtx-pot}}\ge \frac{a+1}{2}\,, \]
		so the minimum mean weight of $(G(S),\w)$ is at least $\frac{a+1}{2}$. Thus, we only need to prove the claim.
		\begin{claimproof}
			Note that since~$u<\frac{a-3+6b}{8}$ is equivalent to $0<\frac{a+1}{2}-(4u-3b+2)$ and~$\frac{3}{4}b\le u$ is equivalent to $\frac{a+1}{2}-(4u-3b+2)\le \frac{a+1}{2}-2$, for all~$u\in S$ we have~$0\le\pi(u)\le \frac{a+1}{2}-1$ and for all~$u,u'\in S$ with~$u\le u'$ we have~$\pi(u)\ge \pi(u')$.
			
			Let~$\ell,\ell'\in S$. If~$\ell=\ell'$, then by \Cref{lem:gap-lower-bound} we have~$\w^{\pi}(\ell,\ell')=\w(\ell,\ell')\ge \ell \ge\frac{a+1}{2}$. If~$\ell>\ell'$, then since~$\ell$ and~$\ell'$ have the same parity by the definition of~$S$, we have~$\ell'<\ell-1$. Furthermore, we have~$\frac{2}{3}\ell\le\frac{2}{3}b<a\le \ell'$. Hence, by \Cref{lem:gap-lower-bound} we have
			\[ \w^\pi(\ell,\ell') = \w(\ell,\ell')+\pi(\ell)-\pi(\ell') \ge \ell'+0-\left(\frac{a+1}{2}-1\right) \ge \frac{a+1}{2}\,.\]
			From now onwards, we have~$\ell<\ell'$. Since~$\ell$ and~$\ell'$ have the same parity, we have~$\ell<\ell'-1$. Furthermore, we have~$\frac{2}{3}\ell'\le\frac{2}{3}b<a\le\ell$. We distinguish two main cases.
			
			\textbf{Case 1:}~$\ell<\frac{3}{4}\ell'$. Since~$\ell<\frac{3}{4}\ell'\le\frac{3}{4}b$, we have~$\pi(\ell)=\frac{a+1}{2}-1$. Since~$b<\frac{3}{2}a$ and~$a\le\ell<\frac{3}{4}\ell'$, we have~$\frac{a-3+6b}{8} <\frac{a+6\cdot\frac{3}{2}a}{8}<\frac{4}{3}a<\ell'$. Hence, we have~$\pi(\ell')=0$. Now by \Cref{lem:gap-lower-bound} we have~$\w(\ell,\ell')\ge1$, so we obtain
			\[ \w^\pi(\ell,\ell') = \w(\ell,\ell') + \pi(\ell) - \pi(\ell') \ge 1 + \left(\frac{a+1}{2} - 1\right) - 0 =\frac{a+1}{2}\,. \]
			
			\textbf{Case 2:}~$\ell\ge\frac{3}{4}\ell'$. In this case, by \Cref{lem:gap-lower-bound} we have~$\w(\ell,\ell')\ge 4\ell-3\ell'+2\ge 2$. We distinguish three subcases.
			
			\textbf{Subcase 2A:}~$\ell<\frac{3}{4}b$. In this subcase, we have~$\pi(\ell)=\frac{a+1}{2}-1$. If~$\frac{a-3+6b}{8}\le\ell'$, then we have~$\pi(\ell')=0$ and so
			\[ \w^\pi(\ell,\ell')=\w(\ell,\ell')+\pi(\ell)-\pi(\ell')\ge 2+\left(\frac{a+1}{2}-1\right)-0\ge\frac{a+1}{2}\,. \]
			If~$\frac{3}{4}b\le\ell'<\frac{a-3+6b}{8}$, then we have~$\pi(\ell')=\frac{a+1}{2}-(4\ell'-3b+2)$ and so
			\begin{align*}
				\w^\pi(\ell,\ell') &= \w(\ell,\ell')+\pi(\ell)-\pi(\ell') \\
				&\ge (4\ell-3\ell'+2)+\left(\frac{a+1}{2}-1\right)-\left(\frac{a+1}{2}-(4\ell'-3b+2)\right)\\
				&=4\ell+\ell'-3b+3\ge5a-\frac{9}{2}a+3\ge\frac{a+1}{2}\,,
			\end{align*}
			where we used~$a\le\ell\le\ell'$ and~$b<\frac{3}{2}a$. Otherwise, we have~$\ell'<\frac{3}{4}b$, in which case we have~$\pi(\ell')=\frac{a+1}{2}-1$ and so
			\[ \w^\pi(\ell,\ell') = \w(\ell,\ell')+\pi(\ell)-\pi(\ell') \ge 4\ell-3\ell'+2 \ge 4a-\frac{27}{8}a +3 \ge\frac{a+1}{2}\,,\]
			where we use~$a\le\ell$, $\ell'<\frac{3}{4}b$ and~$b<\frac{3}{2}a$ to obtain the penultimate inequality.
			
			\textbf{Subcase 2B:}~$\frac{3}{4}b\le\ell<\frac{a-3+6b}{8}$. In this subcase, we have~$\pi(\ell)=\frac{a+1}{2}-(4\ell-3b+2)$. If~$\frac{a-3+6b}{8}\le\ell'$, then we have~$\pi(\ell')=0$ and so
			\begin{align*}
				\w^\pi(\ell,\ell') &= \w(\ell,\ell')+\pi(\ell)-\pi(\ell') \\
				&\ge (4\ell-3\ell'+2)+\left(\frac{a+1}{2}-(4\ell-3b+2)\right)-0 \ge\frac{a+1}{2}\,,
			\end{align*}
			where we simplify expressions and use~$\ell'\le b$ in the final inequality. Otherwise, we have~$\frac{3}{4}b\le\ell'<\frac{a-3+6b}{8}$, in which case we have~$\pi(\ell')=\frac{a+1}{2}-(4\ell'-3b+2)$ and so
			\begin{align*} 
				\w^\pi(\ell,\ell')&= \w(\ell,\ell')+\pi(\ell)-\pi(\ell') \ge 4\ell-3\ell'+2-4\ell+4\ell'=\ell'+2 \ge \frac{a+1}{2}\,, 
			\end{align*}
			where we use~$a\le\ell'$ in the final inequality.
			
			\textbf{Subcase 2C:}~$\frac{a-3+6b}{8}\le\ell$. In this case, we have~$\pi(\ell)=\pi(\ell')=0$. Hence, we have
			\[ \w^\pi(\ell,\ell')= \w(\ell,\ell')+\pi(\ell)-\pi(\ell') \ge 4\ell-3\ell'+2 \ge \frac{a-3+6b}{2}-3\ell'+2 \ge \frac{a+1}{2}\,, \]
			where we used~$\frac{a-3+6b}{8}\le\ell\le\ell'\le b$.
		\end{claimproof}
	\end{proof}

\end{document}
